\section{Introduction}

Unmanned aerial vehicles (UAVs) have become indispensable tools across diverse domains including logistics delivery, power line inspection, forest fire monitoring, and disaster response \cite{icpo2024,qcpo2025,jayacpo2025}. A critical enabler for autonomous UAV operations is three-dimensional (3D) path planning—the task of generating a feasible, safe, and optimal flight trajectory from a start point to a destination while satisfying kinematic constraints and avoiding obstacles \cite{sps2021}. This problem is fundamentally challenging because it requires simultaneously optimizing multiple conflicting objectives under highly nonlinear constraints imposed by terrain, threats, and UAV dynamics.

Existing approaches to UAV path planning can be broadly categorized into graph search methods (e.g., Voronoi diagrams, probabilistic roadmaps), cell decomposition algorithms (e.g., A*, D*), artificial potential field methods, and nature-inspired metaheuristics \cite{sps2021}. Among these, metaheuristic algorithms have emerged as the predominant paradigm due to their gradient-free nature, strong global search capability, and flexibility in handling complex constraints \cite{icpo2024,multiicpo2026}. Particle Swarm Optimization (PSO) and its variants have been particularly popular \cite{sps2021}, alongside more recent algorithms such as Grey Wolf Optimizer (GWO) \cite{gwo2014}, Whale Optimization Algorithm (WOA) \cite{woa2016}, and Crested Porcupine Optimizer (CPO) \cite{cpo2024}.

A landmark advancement in metaheuristic-based UAV path planning was the introduction of \textbf{spherical vector encoding} by Phung and Ha in their Spherical Vector-based PSO (SPSO) \cite{sps2021}. Unlike conventional approaches that directly encode path waypoints as Cartesian coordinates $(x, y, z)$, SPSO represents each path segment as a spherical vector $(\rho, \psi, \phi)$, where the magnitude $\rho$, elevation angle $\psi$, and azimuth angle $\phi$ correspond directly to the UAV's speed, climbing angle, and turning angle, respectively. This encoding offers two decisive advantages: (i) UAV kinematic constraints can be naturally incorporated as search bounds on $\psi$ and $\phi$, dramatically reducing the feasible search space; and (ii) the search operates in the UAV's configuration space rather than Cartesian space, yielding higher-quality solutions. SPSO demonstrated superior performance over multiple PSO variants and other metaheuristics on real Digital Elevation Model (DEM) maps \cite{sps2021}.

The Crested Porcupine Optimizer (CPO), proposed by Abdel-Basset et al. \cite{cpo2024}, is a recently developed metaheuristic inspired by the four defensive behaviors of crested porcupines: visual intimidation, acoustic deterrence, olfactory repellent, and physical attack. These strategies are organized into exploration (visual and acoustic) and exploitation (olfactory and physical) phases, providing CPO with a rich behavioral repertoire. CPO has demonstrated outstanding performance on CEC benchmark suites \cite{cpo2024} and has rapidly attracted improvement efforts. Liu et al. \cite{icpo2024} proposed ICPO, integrating Symbiotic Organisms Search (SOS) mutualism and a periodic retreat strategy, achieving a 63\% cost reduction on urban terrain. Liu et al. \cite{qcpo2025} developed QCPO, embedding Q-learning for adaptive parameter tuning, while Zhang et al. \cite{jayacpo2025} proposed JAYA-CPO, fusing the JAYA algorithm with CPO's attack strategy for power inspection path planning. Zhang et al. \cite{multiicpo2026} introduced Multi-ICPO with Tent chaotic mapping and Cauchy-Gaussian mutation for multi-UAV scenarios. Shi and Li \cite{mocpo2026} extended CPO to multi-objective optimization as MOCPO. Despite these advances, a critical observation emerges: \textbf{none of the existing CPO improvements utilize spherical vector encoding}. All operate in Cartesian space, where the search variables bear no physical relationship to UAV motion parameters, and kinematic constraints must be enforced through penalty terms rather than being intrinsically encoded.

When CPO is naively transplanted into the spherical vector framework, three structural limitations become apparent. \textbf{First}, CPO's visual and acoustic defense strategies rely on random individual differentiation—for instance, the sound strategy computes $\mathbf{x}_{r1} - \mathbf{x}_{r2}$ to generate a search direction. In the spherical coordinate space, where constraint violations produce infinite cost, randomly selected individuals are frequently infeasible, rendering differential guidance ineffective. \textbf{Second}, the exploration–exploitation switch is governed by a purely random mechanism $\text{rand}() < \text{rand}()$, which maintains a fixed 50\% exploration probability throughout the entire optimization process. This prevents the algorithm from naturally transitioning from broad exploration to focused exploitation as the search progresses. \textbf{Third}, CPO lacks a dedicated escape mechanism; once the population converges to a local optimum in the later stages of optimization, none of the four defense strategies can effectively redirect the search.

To address these limitations, this paper proposes \textbf{Sphere-ICPO}, a spherical vector-based improved Crested Porcupine Optimizer for UAV 3D path planning. Building upon the CPO defensive framework and the SPSO spherical encoding infrastructure, three targeted innovations are introduced:

\begin{enumerate}
    \item \textbf{SOS Mutualism for Exploration:} Inspired by ICPO \cite{icpo2024}, the original CPO visual and acoustic defense strategies are replaced with symbiotic organisms search (SOS) mutualism. Unlike random individual differentiation, SOS enables each agent to cooperate with a randomly selected partner and jointly move toward the global best solution, providing robust guidance even when the partner is infeasible. This represents the first application of SOS mutualism in the spherical vector space.

    \item \textbf{Adaptive Exploration Ratio:} The purely random exploration switch $\text{rand}() < \text{rand}()$ is replaced with an adaptive nonlinear decay function $a(t) = 2(0.7(1 - t/T)^{0.5} + 0.3)$, enabling a smooth transition from approximately 100\% exploration in early iterations to approximately 30\% in later iterations. This allows the algorithm to explore broadly when the search space is unknown and exploit precisely when approaching convergence.

    \item \textbf{Stagnation-Triggered Periodic Retreat:} Rather than triggering retreat at fixed intervals as in \cite{icpo2024}, the retreat mechanism is activated only when the global best solution remains unchanged for five consecutive iterations. Retreat employs cosine-geometric back-stepping toward the current best position with a logarithmically decaying step size, providing targeted escape from local optima without disrupting ongoing convergence.
\end{enumerate}

The effectiveness of Sphere-ICPO is validated through comprehensive experiments on four benchmark scenarios constructed from real LiDAR DEM data of Christmas Island, Australia. The threat configurations span from 3 scattered threats to 7 dense threats, covering a complete spectrum of environmental complexity. An ablation study isolates the contribution of each innovation, while multi-algorithm comparisons against SPSO, GWO, CPO, and WOA—supported by Wilcoxon rank-sum and Friedman statistical tests—demonstrate the robustness and statistical significance of the proposed approach. Notably, with an equivalent population of 500 particles, Sphere-ICPO outperforms SPSO on three of four maps while reducing the mean cost by up to 26.7\% relative to the baseline spherical CPO.

The remainder of this paper is organized as follows. Section 2 formulates the UAV path planning problem and the spherical vector encoding framework. Section 3 presents the proposed Sphere-ICPO algorithm in detail. Section 4 reports the experimental results, including ablation studies, multi-map comparisons, statistical tests, and path visualizations. Section 5 concludes the paper and outlines future research directions.
